\chapter{Design and Implementation}
\markboth{Design and Implementation}{Design and Implementation}

\section{UWB}
In swarm drone project, one information is central and will determine the achievement of the shaping.
This information is the \textbf{position}. In our case, we are measuring the relative position between drone. \\ 
To achieve this needed, our professor imposed UWB as localisation sensors. Those sensors are based on electromagnetic signal and needed two modules to interact with each others (Find more information in the bibliographic studies \ref{Bibliographical Studies}.

\subsection{Algorithm choice}
As expose in the bibliographic studies, there are differents algorithms which respond to multiples needs. In this project, we need precise localisation data and no interference, furthermore we have maximun 3 drone in the swarm. \\
With those caracterestics, we choose the DS-TWR location algorithm (\ref{DS-TWR}) and the TDMA timing algorithm (\ref{TDMA}). \\
We have implemented them into the ESP32-DW1000 module.The $\href{https://github.com/M1chigan/Projet-S-IoT-Essaim-de-Drone/blob/main/Code/Code_commun/src/uwb_driver.cpp}{code}$ is using C/C++ as software language. 

\subsection{Calibration}

In general, sensors need a calibration to use them efficiently. For UWB module, we need to establish a distance reference. For this project, the reference is the expected distance between drones: 1.20m.
Afterward, we had to adapt the Antenna Delay parameters which represent signal's time in antenna's PCB. Each module has different time delay. To calibrate, we had to choose a module as a reference, here UWB 1. And then modify Antenna Delay of each drone to make them find 1.20m at 1.20m in real life. 
After this step, all drones are calibrated for each other. 
After doing some tests, we had to modify the reference to 1.30m, this modification is explained in the next part \ref{FtFV}.
\newpage
\subsection{Test}
 During the last year project, our colleagues find out that UWB module should be vertically installed to optimize there precision and reach. So we have repeated the same test  in order to verify there result. 
\subsubsection{Face to Face verticaly}
\label{FtFV}
We have carried out same vertical test following this schema :
\begin{figure}[H]
  \centering
  \FtFVertical
  \caption{Configuration schema of the face to face verticaly}
\end{figure}

\begin{figure}[H]
  \centering
  \begin{minipage}[c]{0.48\textwidth}
    \centering
    \footnotesize
    \setlength{\tabcolsep}{3pt}
    \resizebox{\linewidth}{!}{%
    \begin{tabular}{|l|l|l|}
    \hline
        Real distance (cm) & Measured distance (cm) & Error  \\ \hline
        70 & 67 & 3 \\ \hline
        80 & 78 & 2 \\ \hline
        90 & 89,2 & 0,8 \\ \hline
        100 & 100 & 0 \\ \hline
        110 & 110,6 & -0,6 \\ \hline
        120 & 120 & 0 \\ \hline
        130 & 132,4 & -2,4 \\ \hline
        140 & 143,6 & -3,6 \\ \hline
        150 & 154,2 & -4,2 \\ \hline
        160 & 165,7 & -5,7 \\ \hline
        170 & 176,4 & -6,4 \\ \hline
    \end{tabular}}
    \caption{Measure for face to face verticaly referenced at 120cm}
  \end{minipage}\hfill
  \begin{minipage}[c]{0.48\textwidth}
    \centering
    \footnotesize
    \setlength{\tabcolsep}{3pt}
    \resizebox{\linewidth}{!}{%
    \begin{tabular}{|l|l|l|}
    \hline
        Real distance (cm) & Measured distance (cm) & Error  \\ \hline
        70 & 64 & 6 \\ \hline
        80 & 76,2 & 3,8 \\ \hline
        90 & 85,6 & 4,4 \\ \hline
        100 & 97,4 & 2,6 \\ \hline
        110 & 108,2 & 1,8 \\ \hline
        120 & 119,6 & 0,4 \\ \hline
        130 & 130 & 0 \\ \hline
        140 & 140 & 0 \\ \hline
        150 & 152,7 & -2,7 \\ \hline
        160 & 162 & -2 \\ \hline
        170 & 172,2 & -2,2 \\ \hline
    \end{tabular}}
    \caption{Mesures pour module face à face avec une référence posé à 130cm}
  \end{minipage}
\end{figure}
\noindent As we can observe in those table, the precision is decreasing fast when distance is higher than the reference. So place the reference at 130cm would allow us to get better precision at $\pm$ 10cm around 120cm.
\subsubsection{Face to Face horizontaly}
\begin{figure}[H]
  \centering

  % --- Schéma ---
  \resizebox{0.5\linewidth}{!}{\FtFhorizontale}
  \captionof{figure}{Configuration schema of the face to face horizontaly}

  \vspace{1em}

  % --- Tableau en deux moitiés ---
  \small
  \setlength{\tabcolsep}{4pt}
  \begin{minipage}[t]{0.46\textwidth}
    \vspace{0pt}
    \centering
    \begin{tabular}{|c|c|c|}
      \hline
      Real distance (cm) & Measured distance (cm) & Error \\ \hline
      69,4  & 61    & -8,4 \\ \hline
      70    & 62,2  & -7,8 \\ \hline
      75    & 69,6  & -5,4 \\ \hline
      80    & 73,9  & -6,1 \\ \hline
      82,6  & 76,1  & -6,5 \\ \hline
      85    & 79,9  & -5,1 \\ \hline
      90    & 86,6  & -3,4 \\ \hline
      91    & 87    & -4   \\ \hline
      95    & 92,2  & -2,8 \\ \hline
      100   & 96,5  & -3,5 \\ \hline
      100   & 96,8  & -3,2 \\ \hline
      105   & 102,6 & -2,4 \\ \hline
      107,3 & 106,4 & -0,9 \\ \hline
      110   & 108,9 & -1,1 \\ \hline
      115   & 113,4 & -1,6 \\ \hline
      120,5 & 120,9 & 0,4  \\ \hline
      125   & 126,4 & 1,4  \\ \hline
      130   & 129,6 & -0,4 \\ \hline
      132   & 133,7 & 1,7  \\ \hline
    \end{tabular}
  \end{minipage}\hfill
  \begin{minipage}[t]{0.46\textwidth}
    \vspace{0pt}
    \centering
    \begin{tabular}{|c|c|c|}
      \hline
      Real distance (cm) & Measured distance (cm) & Error \\ \hline
      135   & 133,6 & -1,4 \\ \hline
      140   & 138,6 & -1,4 \\ \hline
      145   & 144,9 & -0,1 \\ \hline
      147,4 & 150   & 2,6  \\ \hline
      150   & 151,7 & 1,7  \\ \hline
      155   & 160,3 & 5,3  \\ \hline
      155,5 & 161   & 5,5  \\ \hline
      160   & 165,7 & 5,7  \\ \hline
      164,7 & 171,8 & 7,1  \\ \hline
      165   & 172,2 & 7,2  \\ \hline
      167   & 172,8 & 5,8  \\ \hline
      170   & 177,2 & 7,2  \\ \hline
      175   & 180,3 & 5,3  \\ \hline
      176   & 183,8 & 7,8  \\ \hline
      180   & 184,3 & 4,3  \\ \hline
      185   & 188,3 & 3,3  \\ \hline
      189,6 & 197,9 & 8,3  \\ \hline
      190   & 191,9 & 1,9  \\ \hline
    \end{tabular}
  \end{minipage}

  \captionof{table}{Mesure face à face de façon horizontale}
\end{figure}

Let's compare both position (horizontal or vertical), as we can observe around 120cm for horizontal position, we got around 1.8cm of error and for vertical, we got around 1cm. Also, when the distance move off 120 cm, the precision decreases slower in vertical than in horizontal. Based on this experimentation, we choose to use our module verticaly to get better precision and flexibility. \\

\noindent Normally, in the documentation, there is an antenna diagram which describe the efficient area of the antenna. We were looking for this diagram in order to choose the shape of the drone which get the best precision. After research, we did not find it, therefore we had to realize some tests in different disposition.

\subsubsection{Paralel module}

\begin{figure}[H]
  \centering
  \begin{minipage}[c]{0.38\textwidth}
    \centering
    \resizebox{\linewidth}{!}{\SchemaParalleles}
    \captionof{figure}{Schema of the paralel configuration}
  \end{minipage}\hfill
  \begin{minipage}[c]{0.58\textwidth}
    \centering
    \footnotesize
    \setlength{\tabcolsep}{3pt}
    \resizebox{\linewidth}{!}{%
    \begin{tabular}{|c|c|c|c|c|}
      \hline
      offset (cm) & Angle (°) & Real distance (cm) & Measured distance (cm) & Error \\ \hline
      0 & 0 & 120 & 117 & 3 \\ \hline
      10 & 4,78 & 121 & 118,4 & 2,6 \\ \hline
      20 & 9,59 & 121,5 & 118,9 & 2,6 \\ \hline
      30 & 14,48 & 122,5 & 119,1 & 3,4 \\ \hline
      40 & 19,47 & 125,8 & 124,8 & 1 \\ \hline
      50 & 24,62 & 128,5 & 128,7 & -0,2 \\ \hline
      60 & 30,00 & 132,5 & 134,4 & -1,9 \\ \hline
      70 & 35,69 & 137 & 142,6 & -5,6 \\ \hline
    \end{tabular}}
    \captionof{table}{Measure of the paralel configuration}
  \end{minipage}
\end{figure}


\subsubsection{Rayon configuration}
\begin{figure}[H]
  \centering
  \begin{minipage}[c]{0.38\textwidth}
    \centering
    \resizebox{\linewidth}{!}{\Schemarayon}
    \captionof{figure}{Schema of the rayon configuration}
  \end{minipage}\hfill
  \begin{minipage}[c]{0.58\textwidth}
    \centering
    \footnotesize
    \setlength{\tabcolsep}{3pt}
    \resizebox{\linewidth}{!}{%
    \begin{tabular}{|l|l|l|}
    \hline
        Angle(°) & Real distance (cm) & Measured distance (cm)  \\ \hline
        10 & 120 & 116,7 \\ \hline
        20 & 120 & 118,4 \\ \hline
        30 & 120 & 119,5 \\ \hline
        40 & 120 & 122,8 \\ \hline
        50 & 120 & 125,9 \\ \hline
        60 & 120 & 124,2 \\ \hline
        70 & 120 & 123 \\ \hline
        80 & 120 & 124 \\ \hline
        90 & 120 & 130 \\ \hline
    \end{tabular}}
    \captionof{table}{Measure of the rayon configuration}
  \end{minipage}
\end{figure}

\subsubsection{Face to Back}
\begin{figure}[H]
  \centering
  \begin{minipage}[c]{0.38\textwidth}
    \centering
    \resizebox{\linewidth}{!}{\FtB}
    \captionof{figure}{Schema of the face to back configuration}
  \end{minipage}\hfill
  \begin{minipage}[c]{0.58\textwidth}
    \centering
    \footnotesize
    \setlength{\tabcolsep}{3pt}
    \resizebox{\linewidth}{!}{%
    \begin{tabular}{|l|l|l|}
    \hline
        Real distance (cm) & Measured distance (cm) & Error  \\ \hline
        70 & 97 & -27 \\ \hline
        80 & 109,8 & -29,8 \\ \hline
        90 & 122,8 & -32,8 \\ \hline
        100 & 133,3 & -33,3 \\ \hline
        110 & 143,5 & -33,5 \\ \hline
        120 & 152 & -32 \\ \hline
        130 & 161,6 & -31,6 \\ \hline
        140 & 177,2 & -37,2 \\ \hline
        150 & 188 & -38 \\ \hline
        160 & 192,2 & -32,2 \\ \hline
        170 & 208,8 & -38,8 \\ \hline
    \end{tabular}}
    \captionof{table}{Measure of the face to back configuration}
  \end{minipage}
\end{figure}

\newpage
\subsection{Filter}
Electromagnetics signals are sensitive to exterior noise and can be interfered easily. Therefore, we have to filter the raw data. To do so we implement three differents filters.  
\subsubsection{Polynomial filter}
During calibration and test, we have discovered that the errors is linked to the real distance and variate following a polynomial equation. 

\begin{figure}[H]
   \centering
   \resizebox{0.6\textwidth}{!}{\includegraphics{Polynomial.png}}
   \caption{Diagram of the error based on real distance}
\end{figure}

\noindent Using the Excel linearization tool, we got this error equation: 

\begin{equation}
   f(x) = -1{,}004 \times 10^{-6}\, x^4 + 0{,}0005059\, x^3 - 0{,}09263\, x^2 + 7{,}398\, x - 221{,}7
\end{equation}
Every measure pass through the equation and are recalibrate based on this equation. That permit to reduce the error when we are at more $\pm$ 10 cm from the reference.\\

\noindent We have implemented it in C/C++ : 
\begin{figure}[H]
   \centering
   \resizebox{0.7\textwidth}{!}{\includegraphics{Polynomial_Filter.png}}
   \caption{\href{https://github.com/M1chigan/Projet-S-IoT-Essaim-de-Drone/blob/main/Code/Code_commun/src/filtre.cpp}{Code} for polynomial filter}
\end{figure}
\newpage
\subsubsection{NLOS Filter}
Electromagnetic signal are very sensitive to obstacle. After doing research, we have discovered that the signal will lost power if it encounters an obstacle. Based on that, we can verify that the power receives is strong enough otherwise the value cannot be considered as feasible so we can dennied it.\\

\noindent We have implemented it in C/C++: 

\begin{figure}[H]
   \centering
   \resizebox{0.70\textwidth}{!}{\includegraphics{NLOS Filter.png}}
   \caption{\href{https://github.com/M1chigan/Projet-S-IoT-Essaim-de-Drone/blob/main/Code/Code_commun/src/filtre.cpp}{Code} for NLOS Filter}
\end{figure}


\subsubsection{Innovation Gate Filter}
UWB has tendencies to be noisy, that create some absurd value very high or very low. This data can affect the corrector and create some discontinuity in the drone control. Based on the max speed of the drone, we can determine absurdity by calculating the difference $d_{-1}$ and $d$, then determine if the difference can be reach at the max speed, if not the measure is ignored, otherwise, we accepted it.\\

\noindent We have implemented it in C/C+: 

\begin{figure}[H]
   \centering
   \resizebox{0.70\textwidth}{!}{\includegraphics{Innovation_Gate.png}}
   \caption{\href{https://github.com/M1chigan/Projet-S-IoT-Essaim-de-Drone/blob/main/Code/Code_commun/src/filtre.cpp}{Code} for innovation gate filter}
\end{figure}

\newpage

\section{Communication}
In the project, we have to differents use for communication, first communication is used to made each drone talk with each other to maintain the shape but also adapt themselves to the other. Second communication is the data's descent from the drone to the base to permit visualition of the swarm and command drone.

\subsection{Choice}
We choose ESP-NOW has the protocol for two main raison : it doesn't need new hardware implementation. And, it responds to both communication purpose of the project. Then, has said during the presentation, the code has been optimized and write by the manufacturer which facilitate it implementation. For exemple, send data is only this command :
\begin{figure}[H]
   \centering
   \resizebox{0.7\textwidth}{!}{\includegraphics{ESP-NOW_send.png}}
   \caption{Code for send messsage using ESP-NOW}
\end{figure}

Has said before, range reach 150 meters allows us to create a base. It only ask  us, to have a ESP32, connect to our computer.

\subsection{ESP-NOW implementation}
To implement ESP-NOW, we had to think about the data we want to exchange. For our first test, we are using two different packages for two different purpose. \\
First for inter-drone communication : 
\begin{figure}[H]
   \centering
   \resizebox{0.6\textwidth}{!}{\includegraphics{Drone_package.png}}
   \caption{Drone communication package}
\end{figure}

And for the base to drone communcation we have this package : 

\begin{figure}[H]
   \centering
   \resizebox{0.6\textwidth}{!}{\includegraphics{base_package.png}}
   \caption{Drone communication package}
\end{figure}


\newpage

\section{Measure processing and formation control}

After studying the UWB measurements and the communication with ESP-NOW, we have to consider the processing of the retrieved data to finally implement an optimal formation control algorithm. The questions that need to be asked are the following : what data do we need to retrieve and what algorithm do we use for its processing and for the formation control. As mentioned before, our drones include many onboard sensors and benefit from the added UWB module. The bibliographical study gave us insights about potential architectures of measure and control so a choice has to be made considering the hardware available. Here are two things to notice :


\begin{itemize}
    \item UWB modules return purely relative distance measurements and give no sense of direction at all
    \item The onboard Extended Kalman Filter from the flight controller's firmware Ardupilot fuses the measurements from its different sensors (accelerometer, gyroscope, magnetometer and barometer) and the added GNSS module to get notably an accurate estimation of the drone's attitude (yaw, roll, pitch) as a quaternion, 3 NED (North East Down) velocity components and 3 NED position components (up to 22 states are estimated in total).
    
\end{itemize}

 Therefore, we know that we can get good estimates of these states through the onboard EKF of the flight controller. A first choice has to be made about what we want to control. 
\newpage

\section{Visualization}
In outdoor test, drones are flying at more than 10 meters away from us. That constraint reduces our visibility and comprehension of the drone movements, attitude and position. To prevent this problematic we need to develop a visualization system. That would increase our error detection capability and give the opportunity to save test and simulate them afterward to analyse and detect problematic point.\\
Visualization would also allow us to observe in real time differents parameters as: attitude, 3D shape positionning, and also RAW data coming from the drone.\\

In order to achieve this system, we choose the sofware Matlab as a visualizer and software language.  

\subsection{Matlab}

Matlab is a general tool developped by the MathWorks. It has many functions such as mathematics calculator, simulator, and also a corrector designer. Based on add-on and toolbox, it can be usefull for many differents purpose. It also permits to generate an interface. 

\subsection{Architecture}

The visualization follow the next architecture : 
\begin{figure}[H]
  \centering
   \resizebox{0.9\textwidth}{!}{\SchVisualisation}
  \caption{Visualization architecture}
\end{figure}

Let's explain each different bloc and its role : 

\begin{itemize}
    \item ESP32 : Relay between ESP-NOW protocol and the UART Port of the computer, which permit to make it visible for the matlab code.
    \item COM PORT: Physical bridge between the computer and the ESP32
    \item mainMenu : General window which redirect for the next window
    \item runObserver : Window to observe UART communication on the PORT (Debug utility)
    \item runSimulation : Main Window where we visualise attitude of the drone, shape placement and data reader
\end{itemize}

\newpage
\section{Simulation}
\label{Simulation}

Due to logistical problem, outdoor test are rare, and INSA Strasborug does not have indoor test room. Howewer our work need to be tested and confronte reality. To achieve this purpose, we had to develop a simulation protocol which is the nearest possible from the reality. \\
To do so, we had to choose same firmware as our flight controler and an UAV simulation tool. \\
In this part, we will first present both firmware : ArduPilot and Mission planer, second, explain the simulation architecture, then present the result we obtained and finally the simulation possibilities.

\subsection{ArduPilot and Mission Planer}

To do drone simulation, we have to select the same architecture as our drone on the pixhwak 4. In the drone industry, two differents architectures are existing: Ardupilot and PX4. In our case, pixhwak 4 is using the ArduPilot architectures. \\
Then, after selecting the architecture, we have to use a software to visualise and simulate drone. We choose mission planer because we have already work with it on different project. \\
Let's present both tools:

\subsubsection{\href{https://github.com/ArduPilot/ardupilot}{ArduPilot}}

 ArduPilot was created in 2007 by the DIY Drones community. Originally developed for Arduino boards, it has since become a general-purpose, open-source autopilot. It can control many types of vehicles (copters, planes, rovers, submarines, etc.). ArduPilot provides autopilot functions, receives data from sensors, and fuses it using an Extended Kalman Filter (EKF) [ref]. It also allows users to create pre-planned missions. To operate, it uses the MAVLink protocol to communicate with the drone and facilitate control. In our case, an ESP32 uses the MAVLink library to send MAVLink messages that command the drone’s movements.

\subsubsection{\href{https://ardupilot.org/planner/}{Mission planer}}

Mission planer is a representative tool for copters such as drone. Created by  Michael Oborne to help users of ArduPilot architecture to vizualise their drone and command them. The software also permits to declare  multiple drone so create a swarm. For our project, it will be used as a simulation field. 
\newpage
\subsection{Architecture}

The simulation works with a very specific architecture cause tools like ArduPilot works under Linux system. The architecture is exposed in the figure below. 

\begin{figure}[H]
  \centering
   \resizebox{0.9\textwidth}{!}{\SchSimulation}
  \caption{Simulation architecture}
\end{figure}

Let's explain each different bloc and its role : 
\begin{itemize}
    \item $ESP-32_{Base}$ : Represent order send from the base for exemple takeoff
    \item $ESP-32_{Drone}$: Represent the main drone where code we want to test is running
    \item ArduPilot : Represent the pixhawk of our drone. Here, we are simulating 3 of them. 
    \item \href{https://github.com/M1chigan/Projet-S-IoT-Essaim-de-Drone/blob/main/Simulation/Script%20d'essai/Bridge.py}{Bridge.py} : Code made to redirect data from the UART connection to UDP connection 
    \item MAVProxy : Manage all entries and exit permit to link the drone entries redirect them to the Ardupilot and also permit to redirect exit from ArduPilot (GPS, speed, etc) to Mission Planer.
    \item Mission Planer : Use to visualise our swarm drone based on data from ArduPilot
    \item \href{https://github.com/M1chigan/Projet-S-IoT-Essaim-de-Drone/blob/main/Simulation/Script%20d'essai/Simulation.py}{Simulation.py}  : Python scrip to simulate $Drone_2$ and $Drone_3$
\end{itemize}


\newpage
\section{Test}
Test is the main part of technical project, reality is always different from the theory. Therefore, we have to carry out them through two different environement: simulation and outdoor.
\subsection{Simulation}
Using the architecture describe in the part \ref{Simulation}, we did and planned to do different test for different axes of the project.  
\subsubsection{Ongoing : }
\label{Simu-level1}
\underline{Console control:} \\
This test has for purpose to test MAVLink function, communication and drone control.\\ 
\noindent It follows the next schema: \\

\begin{figure}[H]
   \centering
   \resizebox{0.7\linewidth}{!}{\LevelUN}
   \caption{Architecture of the first level of simulation}
\end{figure}

\noindent The drone is for now capable of this function (lexicon below) : 

\begin{figure}[H]
   \centering
   \resizebox{0.3\textwidth}{!}{\includegraphics{lexique.png}}
   \caption{Drone capability at level 1}
\end{figure}

\noindent We can conclude that our communication between base and drone, is well-designed and the MAVLink procedure is well send to the drone. 

\subsubsection{Planned :}

\noindent  \underline{Level 2:} \\ 
\label{Simu-level2}
Objective of this level, a slave drone need to follow the master drone controlled by the console as in the level 1. To do so it should use communication, attitude and GPS data. Possibly, we can think of integrate UWB to make a full test, closest to the real situation.   

\begin{figure}[H]
   \centering
   \resizebox{0.8\linewidth}{!}{\LevelDeux}
   \caption{Architecture of the second level of simulation}
\end{figure}

\noindent \underline{Level 3:} \\ 
Objective of this level is to reach 3 drone and introduce swarm drone. On the 3 drones, two of them would be control by a python script. The third drone is the tested one. It permits to test corrector and attitude corrector. 

\begin{figure}[H]
   \centering
   \resizebox{0.8\linewidth}{!}{\LevelTrois}
   \caption{Architecture of the third level of simulation}
\end{figure}

\newpage
\noindent \underline{Level 4:} \\
Objective of this level is to reach a full simulation environement where all the three drone are working on it own ESP32. Mission planer and Ardupilot are strictly an environement simulation. It is the closest simulation to the objective of the project.

\begin{figure}[H]
   \centering
   \resizebox{0.8\linewidth}{!}{\LevelQuatre}
   \caption{Architecture of the fourth level of simulation}
\end{figure}

\subsection{Outdoor}
Make drone flying is restricted by the laws. It makes test difficult to proceed and rare. Therefore, outdoor test needs to be precisely prepare to not lose any time and be efficient on site.
\subsubsection{Planned}
We have planned to attempt outdoor test within 2 weeks. During this test, we would love to test: 
\begin{itemize}
  \item UWB placement (rotation, and heigh)
  \item ESP-NOW range 
  \item Drone control console (as simulation level 1 [\ref{Simu-level1}])
  \item Slave drone following master drone (as simulation level 1 [\ref{Simu-level2}])
\end{itemize} 


